\section{When to Use Merge-Split}
\label{sec:when-to-use}

\subsection{Decision Framework}

The choice among standard \recom{}, \mergesplit{}, \forestrecom{}, and
SMC~\citep{dellaLibera2026smc} depends on three factors:

\begin{enumerate}
  \item \textbf{Distributional requirement.} Does the application require
        exact targeting of the uniform distribution, or is approximate
        uniformity sufficient?
  \item \textbf{Computational budget.} How many steps (and wall-clock
        time) are available?
  \item \textbf{Ensemble size.} Is the goal to produce a large ensemble
        ($\geq 10{,}000$ plans) or a small focused ensemble
        ($\leq 1{,}000$ plans)?
\end{enumerate}

\subsection{Use Merge-Split When}

\mergesplit{} is the recommended chain when:

\begin{itemize}
  \item Approximate uniformity is sufficient --- the application benefits
        from a known (approximately uniform) stationary distribution, but
        does not require formal exact-determinant guarantees.
        This covers most expert-witness and legislative-disclosure contexts
        where the practitioner can state ``the chain approximately targets
        the uniform distribution'' without requiring mathematical exactness.

  \item Speed matters --- the ensemble must be large (many thousands of
        plans) for statistical power, or the state is large ($n \gtrsim
        2{,}000$ tracts), making \forestrecom{}'s $O(m^2)$ cost per step
        impractical within the available budget.

  \item Standard ReCom's distributional ambiguity is a concern --- when
        opposing counsel or a court scrutinises the ensemble methodology,
        the unknown stationary distribution of standard \recom{} is a
        vulnerability.
        \mergesplit{} provides a principled, citable response~\citep{janson2022}.

  \item The acceptance rate of \forestrecom{} ($\approx\!40\%$) is too
        low for the available computational budget --- \mergesplit{}'s
        $\approx\!60\%$ acceptance rate produces~50\% more accepted plans
        per step, which may be decisive for small-budget applications.
\end{itemize}

\subsection{Use Forest ReCom When}

\forestrecom{}~\citep{dellaLibera2026forestRecom} is preferred when:

\begin{itemize}
  \item Strict theoretical guarantees are required --- the application
        must demonstrate exact detailed balance rather than approximate
        balance.
        This matters for peer-reviewed publication claims about exact
        uniformity.

  \item The merged regions are small ($m \lesssim 100$ edges) ---
        in this regime, the $O(m^2)$ cost is negligible and the exact
        determinant provides additional precision without cost.

  \item The ensemble is small and variance in the ratio estimate is a
        concern --- with fewer than 100 steps, the random estimate in
        \mergesplit{} may introduce systematic bias in the seat
        distribution, whereas the exact determinant of \forestrecom{} is
        immune to this.
\end{itemize}

\subsection{Use SMC When}

Sequential Monte Carlo~\citep{dellaLibera2026smc} is the gold standard:

\begin{itemize}
  \item Exact calibration is required --- SMC produces samples from the
        exact target distribution without the approximation of
        \mergesplit{} or the forest-measure approximation of
        \forestrecom{}.

  \item The ensemble must provably represent the full plan space ---
        SMC's particle filter guarantees coverage that MCMC methods
        cannot match for fixed step counts.

  \item The chain mixing time is a concern --- SMC does not mix; it
        directly samples from the target distribution.
\end{itemize}
The price of SMC is implementation complexity and memory: retaining and
resampling a particle population requires $O(N \cdot n)$ memory for $N$
particles.
For production 50-state runs, \mergesplit{} is usually preferred.

\subsection{Practical Guidance}

For standard 50-state redistricting analysis:

\begin{enumerate}
  \item Generate the initial plan using \cs{}~\citep{b11paper} or
        another global search method.
  \item Run \mergesplit{} for 1{,}000--10{,}000 steps to build the
        ensemble.
  \item For litigation contexts, state explicitly that the chain uses the
        Merge-Split proposal~\citep{janson2022} with approximate
        uniformity targeting.
  \item If the state has large districts ($n/k \geq 200$ tracts per
        district), prefer a longer run to compensate for the larger
        merged-region size.
  \item Record the base seed, step count, and the best plan's district
        assignment hash for reproducibility.
\end{enumerate}

\paragraph{CLI reference.}
\begin{verbatim}
# 1000-step Merge-Split chain, best plan (percentile 0.0)
redist state --state NC --year 2020 \
    --search merge-split --merge-split-steps 1000 --percentile 0.0

# 5000-step chain for a larger ensemble
redist state --state TX --year 2020 \
    --search merge-split --merge-split-steps 5000 --percentile 0.0
\end{verbatim}
